% ==============================================================================
% ssec:cs_penny --- Případová studie: PKS Penny Market 2026
% Zdroje: PennyMarket_KS2026, CMKOS_ZpravaKV2025, CMKOS_PrinosKV2025,
%         zakon_zp_2006, zakon_kv_1991, Armstrong_HRM2006, smernice_pay_transparency_2023,
%         penny_top_employer_2026
% ==============================================================================

\Acf{PKS} ve společnosti \emph{Penny~Market, s.\,r.\,o.} na~rok 2026 představuje příklad podnikového vyjednávání v~českém maloobchodě (\acs{NACE}~47.11) s~nízkou odborovou organizovaností a~ročním vyjednávacím cyklem. Smlouva byla uzavřena 20.~února 2026 v~Radonicích bez~automatické prolongace. Rozbor sleduje motivace a~podmínky, jimiž obě strany vyjednávání podmiňují a~co jím získávají.\par

Na~straně zaměstnanců vystupuje koalice pěti odborových subjektů: Odborový svaz pracovníků zemědělství a~výživy \acs{OSPZV}--\acs{ASO}~\acgen{ČR} jako zastřešující svaz a~čtyři základní organizace. Koordinace pod celostátním svazem na~základě plné moci z~roku 2013 předchází potenciální vyjednávací blokádě, přičemž zastřešujícím partnerem v~maloobchodě se~stává svaz z~příbuzného sektoru (zemědělství a~potravinářství).\par

Zastřešující svaz je~přímou smluvní stranou této \aclgen{PKS} podle \mbox{§~25} odst.~1 \aclgen{ZP}. Smluvní standardy tak platí z~titulu podnikového vyjednávání, nikoli na~základě odvětvové extenze. Souběžná \acs{KSVS} pro~zemědělství a~potravinářství, rozšířená pro rok 2026 dle~\mbox{§~7} \aclgen{ZKV} (kap.~\ref{par:cs_ksvs_zemedelstvi}), se~týká pouze odvětví \acgen{NACE}~01.1 až~01.6, a~na~maloobchodního zaměstnavatele (NACE~47.11) se~nevztahuje.\par

Hlavní protihodnotou pro~zaměstnavatele je~flexibilita organizace pracovní doby. Její základní délka kopíruje zákonné limity (\SI{40}{\hour\per\week} v~jednosměnném režimu bez~zkrácení). Pro prodejny je~však sjednáno vyrovnávací období pro nerovnoměrné rozvržení v~maximální délce \SI{52}{\week}, což je~institut podmíněný právě existencí kolektivní smlouvy. Ostatní parametry jako příplatek za~přesčasy (\SI{25}{\percent}), možnost sjednat mzdu s~přihlédnutím k~přesčasům u~vedoucích (do~\SI{150}{\hour\per\rok}) i~příplatky za~noční a~víkendovou práci zůstávají na~zákonném minimu.~\cite{zakon_zp_2006}\par

Mzdový nárůst pro rok 2026 byl sjednán diferencovaně s~účinností od~1.~března: v~průměru \SI{5}{\percent} pro prodejny a~sklady, \SI{3}{\percent} pro ostatní pozice. Tyto hodnoty zaostávají za~typickým růstem v~podnikových smlouvách, který se podle zprávy \aclgen{ČMKOS} pohybuje kolem \SI{6}{\percent} až~\SI{7}{\percent}. Průměrová formulace navýšení navíc dává zaměstnavateli možnost individuální diferenciace a~zajišťuje mu~vysokou předvídatelnost mzdových nákladů.\par

Kompenzací za~nízký mzdový růst jsou nemzdové benefity: dovolená v~rozsahu \SI{6}{\week} (o~dva týdny nad~zákonný základ), příspěvek na~penzijní připojištění \SI{800}{\czk} po~roce zaměstnání, stravenka v~nominální hodnotě \SI{110}{\czk} (s~příspěvkem firmy \SI{60}{\czk}) a~vánoční dárková karta v~hodnotě \SI{3000}{\czk}. V~rámci osobních překážek je sjednáno zvýhodnění nad zákonný rámec, např.~4~dny volna při vlastní svatbě a~5~dní při úmrtí partnera či dítěte. Tato struktura odpovídá zjištění \aclgen{ČMKOS}, že čeští zaměstnavatelé v~\acp{PKS} upřednostňují nákladově stabilní benefity před mzdovým nadstandardem.~\cite{CMKOS_ZpravaKV2025}\par

Celková bilance vyznívá symetricky: zaměstnavatel získává sociální smír (čl.~II.C.7 vylučuje stávkový střet o~věcech upravených smlouvou) a~maximální časovou flexibilitu vyrovnávacího období, zatímco zaměstnanci získávají nadstandard v~podobě dovolené a~benefitů bez~mzdových nadstaveb. Témata \acs{BOZP} (vymezená podrobnými závazky k~OOPP a~kontrolám) a~vzdělávání, které je sjednáno v~konzultačním režimu bez~konkrétních nároků na~rozvoj kvalifikace.~\cite{PennyMarket_KS2026}\par

Tato dělba se~promítá i~do~vnější komunikace zaměstnavatele. Prezentace značky navenek (např.~certifikace Top~Employer 2026) nezmiňuje existenci \aclgen{PKS} či roli sociálních partnerů, což odpovídá české neviditelnosti kolektivního vyjednávání v~rámci strategického řízení lidských zdrojů a~značky zaměstnavatele. Z~hlediska teorie \acgen{HRM} však uplatnění vyjednaných, transparentních pravidel plní významnou koordinační a~legitimační funkci, kterou nelze nahradit jednostranným deklarováním benefitů ze~strany vedení.~\cite{penny_top_employer_2026}\par
